\makeatletter
\@ifundefined{keepfig}{\newenvironment{keepfig}{\par\edef\keepfig@pi{\the\parindent}\noindent
  \begin{minipage}{\linewidth}\parindent=\keepfig@pi\relax}{\end{minipage}\par}}{}
\makeatother
\section{Сепарабельные, нормальные расширения и расширения Галуа}

\begin{definition}[3.6]
Алгебраическое расширение $E/F$ называется \emph{сепарабельным}, если минимальный многочлен над $F$ каждого элемента
$E$ сепарабелен; в противном случае оно называется \emph{несепарабельным}.
\end{definition}

Каждый унитарный неприводимый многочлен из $F[X]$, имеющий корень $\alpha\in E$, есть минимальный многочлен $\alpha$.
Поэтому $E/F$ сепарабельно тогда и только тогда, когда сепарабелен каждый неприводимый многочлен из $F[X]$, имеющий
корень в $E$. По \rf{2.20} и \rf{2.24} расширение $E/F$ несепарабельно тогда и только тогда, когда
\begin{itemize}
\item поле $F$ несовершенно (в частности, $\operatorname{char}F=p\ne0$; более того, $F$ не алгебраично над $\F_p$,
\rf{2.25}(b)), и
\item найдётся $\alpha\in E$, минимальный многочлен которого имеет вид $g(X^p)$, $g\in F[X]$.
\end{itemize}
В частности, каждое алгебраическое расширение совершенного поля (поля характеристики $0$, конечного поля)
сепарабельно (\rf{2.24a}). Пример несепарабельного расширения: $\F_p(T)/\F_p(T^p)$. Элемент $T^p$ трансцендентен над
$\F_p$, так что $\F_p(T^p)$~--- поле рациональных функций от $T^p$, и $T^p$ не является в нём $p$-й степенью
(\rf{2.25}(d)). Поэтому $X^p-T^p$ неприводим над $\F_p(T^p)$ (\rf{2.3}) и является минимальным многочленом $T$; он
несепарабелен: $X^p-T^p=(X-T)^p$.

\begin{remark}[3.6a]
Если $E/F$ сепарабельно и $F\subset M\subset E$, то $M/F$ и $E/M$ сепарабельны. Для $M/F$ это очевидно. Для $E/M$: если
$\alpha\in E$, то минимальный многочлен $\mu^M_\alpha$ элемента $\alpha$ над $M$ делит $\mu^F_\alpha$ в $M[X]$ (так как
$\mu^F_\alpha\in M[X]$ и $\mu^F_\alpha(\alpha)=0$), а делитель многочлена без кратных корней сам не имеет кратных корней.
Итак, \emph{сепарабельность элемента сохраняется при расширении основного поля}. Обратное утверждение~--- если
$M/F$ и $E/M$ сепарабельны, то $E/F$ сепарабельно,~--- тоже верно, но менее очевидно (\rf{3.15c}).
\end{remark}

\begin{definition}[3.7]
Алгебраическое расширение $E/F$ называется \emph{нормальным}, если минимальный многочлен над $F$ каждого элемента $E$
раскладывается на линейные множители в $E[X]$.
\end{definition}

Иначе говоря, $E/F$ нормально, если каждый неприводимый многочлен из $F[X]$, имеющий хотя бы один корень в $E$,
раскладывается в $E[X]$ на линейные множители. Пусть $E/F$ алгебраично и $f\in F[X]$~--- унитарный неприводимый многочлен
с корнем в $E$ (т.\,е.\ минимальный многочлен элемента $E$). Тогда
\[
\left.\begin{aligned}
E/F\ \text{нормально}\ &\Rightarrow\ f\ \text{расщепляется в}\ E\\
E/F\ \text{сепарабельно}\ &\Rightarrow\ \text{корни}\ f\ \text{простые}
\end{aligned}\right\}\ \Rightarrow\ \text{у}\ f\ \text{ровно}\ \deg f\ \text{различных корней в}\ E.
\]
Следовательно, \emph{$E/F$ нормально и сепарабельно тогда и только тогда, когда каждый неприводимый многочлен из
$F[X]$, имеющий корень в $E$, имеет в $E$ ровно $\deg f$ различных корней}.

\begin{example}[3.8]
(a) Многочлен $X^3-2$ имеет один вещественный корень $\sqrt[3]2$ и два невещественных корня в $\CC$
(\rf{3.3}(1)). Поэтому расширение $\Q[\sqrt[3]2]/\Q$ (сепарабельное, так как $\Q$ совершенно) не нормально.

(b) Расширение $\F_p(T)/\F_p(T^p)$ не сепарабельно: минимальный многочлен $X^p-T^p$ элемента $T$ несепарабелен. При этом
оно нормально: это поле разложения $X^p-T^p=(X-T)^p$ (\rf{3.10a} ниже).
\end{example}

\begin{definition}[3.9]
Расширение полей $E/F$ называется \emph{расширением Галуа}, если оно конечно, нормально и сепарабельно. В этом случае
группа $\Aut(E/F)$ называется \emph{группой Галуа} расширения $E/F$ и обозначается $\Gal(E/F)$.
\end{definition}

\begin{theorem}[3.10]
Для расширения $E/F$ следующие условия равносильны:
\begin{enumerate}[label=\textup{(\alph*)}]
\item $E$~--- поле разложения некоторого сепарабельного многочлена $f\in F[X]$;
\item $E$ конечно над $F$ и $F=E^{\Aut(E/F)}$;
\item $F=E^G$ для некоторой конечной группы $G$ автоморфизмов поля $E$;
\item $E/F$~--- расширение Галуа.
\end{enumerate}
\end{theorem}

\begin{keepfig}
\begin{center}
\begin{tikzpicture}[every node/.style={font=\small}]
\node (a) at (0,1.7) {(a) $E$~--- поле разложения сепарабельного $f$};
\node (b) at (8.2,1.7) {(b) $F=E^{\Aut(E/F)}$};
\node (c) at (8.2,0) {(c) $F=E^G$, $|G|<\infty$};
\node (d) at (0,0) {(d) конечно, нормально, сепарабельно};
\draw[cnew,double,-implies] (a) -- node[above]{\rf{3.2}, \rf{3.3c}} (b);
\draw[cnew,double,-implies] (b) -- node[right]{\rf{2.14}} (c);
\draw[cnew,double,-implies] (c) -- node[above]{\rf{3.4}, орбиты} (d);
\draw[cnew,double,-implies] (d) -- node[left]{различные $\mu_{\alpha_i}$} (a);
\end{tikzpicture}
\end{center}
\figcap{Схема доказательства \rf{3.10} и то, на что опирается каждый шаг.}
\end{keepfig}

\begin{proof}
(a)$\Rightarrow$(b). Поле $E$ порождено конечным числом алгебраических элементов (корней $f$) и потому конечно над $F$
(\rf{1.30}). Положим $F'=E^{\Aut(E/F)}$; ясно, что $F\subset F'$, и нужно доказать $F'=F$. Поле $E$ является полем
разложения $f$ и как многочлена из $F'[X]$: оно порождено над $F'$ корнями $f$ (так как порождено ими уже над $F$) и
расщепляет $f$. Многочлен $f$ остаётся сепарабельным над $F'$ (пункт (d) после \rf{2.22}). По \rf{3.2}, применённому
дважды,
\[
|\Aut(E/F')|=[E:F']\le[E:F]=|\Aut(E/F)|.
\]
Но $\Aut(E/F')=\Aut(E/F)$: включение $\subset$ верно, так как $F\subset F'$, а $\supset$~--- по определению $F'$ (каждый
элемент $\Aut(E/F)$ неподвижно оставляет $F'$). Это частный случай \rf{3.3c}(iii); можно сослаться и на \rf{3.5} для
$G=\Aut(E/F)$. Значит, $[E:F']=[E:F]$, откуда $[F':F]=1$ (\rf{1.20}) и $F'=F$.

(b)$\Rightarrow$(c). Положим $G=\Aut(E/F)$. По условию $F=E^G$, а $G$ конечна, так как $E/F$ конечно: $|G|\le[E:F]$ по \rf{2.14}.

(c)$\Rightarrow$(d). По теореме \rf{3.4} $[E:F]\le(G:1)$; в частности, $E/F$ конечно. Пусть $\alpha\in E$ и $f$~--- его
минимальный многочлен над $F$. Нужно показать, что $f$ раскладывается в $E[X]$ на \emph{различные} линейные множители.
Пусть $\{\alpha_1=\alpha,\alpha_2,\dots,\alpha_m\}$~--- орбита $\alpha$ при действии $G$ на $E$ (элементы $\alpha_i$ попарно
различны), и
\[
g(X)=\prod_{i=1}^m(X-\alpha_i)=X^m+a_1X^{m-1}+\dots+a_m.
\]
Каждый $\sigma\in G$ переставляет орбиту ($\sigma(\tau\alpha)=(\sigma\tau)\alpha$, и $\sigma$ инъективен), а коэффициенты $a_j$~---
симметрические многочлены от $\alpha_i$; поэтому $\sigma a_j=a_j$ для всех $\sigma\in G$, т.\,е.\ $g\in F[X]$. Так как $g$
унитарен и $g(\alpha)=0$, он делится на $f$. С другой стороны, если $\alpha_i=\sigma\alpha$, то, применяя $\sigma$ к равенству
$f(\alpha)=0$ (коэффициенты $f$ лежат в $F$ и неподвижны), получаем $f(\alpha_i)=0$. Значит, $f$ делится на попарно
различные (и потому попарно взаимно простые) линейные многочлены $X-\alpha_i$, а значит, и на их произведение $g$. Итак,
$f=g$, и $f$ раскладывается в $E[X]$ на различные линейные множители. Следовательно, $E/F$ нормально и сепарабельно.

(d)$\Rightarrow$(a). Поле $E$ конечно над $F$, поэтому $E=F[\alpha_1,\dots,\alpha_m]$ (например, для базиса
$\alpha_1,\dots,\alpha_m$). Пусть $f_i$~--- минимальный многочлен $\alpha_i$ над $F$, и $f$~--- произведение \emph{различных}
многочленов среди $f_1,\dots,f_m$ (каждый берётся один раз). Так как $E/F$ нормально, каждый $f_i$ раскладывается в
$E[X]$, т.\,е.\ $f$ расщепляется в $E$; поле $E$ порождено корнями $f$ и потому является его полем разложения. Так как
$E/F$ сепарабельно, все $f_i$ сепарабельны, и $f$~--- произведение попарно различных унитарных неприводимых
сепарабельных многочленов~--- сепарабелен (пункт (c) после \rf{2.22}).
\end{proof}

В шаге (d)$\Rightarrow$(a) нельзя брать произведение \emph{всех} $f_i$: минимальные многочлены разных образующих могут
совпадать, даже если ни одна образующая не лежит в поле, порождённом предыдущими. Например, $E=\Q[\sqrt[3]2,\omega\sqrt[3]2]$
($\omega=e^{2\pi i/3}$), $\alpha_1=\sqrt[3]2\in\R$, $\alpha_2=\omega\sqrt[3]2\notin\Q[\alpha_1]$; оба имеют минимальный
многочлен $X^3-2$, и $f_1f_2=(X^3-2)^2$ не сепарабелен.

Сепарабельность в (a) нужна лишь для сепарабельности $E/F$; нормальность от неё не зависит.

\begin{proposition}[3.10a]
Конечное расширение $E/F$ нормально тогда и только тогда, когда $E$~--- поле разложения некоторого многочлена
$f\in F[X]$ (не обязательно сепарабельного).
\end{proposition}

\begin{proof}
$\Rightarrow$. Как в (d)$\Rightarrow$(a): $E=F[\alpha_1,\dots,\alpha_m]$, и $E$~--- поле разложения произведения минимальных
многочленов $\alpha_i$.

$\Leftarrow$. Пусть $E$~--- поле разложения $f$, $f=c\prod_{i=1}^n(X-\alpha_i)$ в $E[X]$, $E=F[\alpha_1,\dots,\alpha_n]$. Пусть
$\beta\in E$, $g$~--- его минимальный многочлен над $F$, и $\Om$~--- поле разложения $g$ над $E$ (\rf{2.7}); тогда
$\Om\supset E$ и $g$ раскладывается в $\Om[X]$. Пусть $\beta'\in\Om$~--- любой корень $g$; покажем, что $\beta'\in E$. По
\rf{2.1}(b) существует $F$-гомоморфизм $\psi\colon F[\beta]\to\Om$ с $\psi(\beta)=\beta'$. Поле $E$ порождено над $F[\beta]$
корнями $f$, а многочлен $\psi f=f$ расщепляется в $\Om$ (уже в $E$); по \rf{2.12a} $\psi$ продолжается до гомоморфизма
$\sigma\colon E\to\Om$. Каждый $\sigma(\alpha_i)$~--- корень $f$ в $\Om$, а корни $f$ в $\Om$~--- это в точности $\alpha_1,\dots,\alpha_n$
(разложение $f=c\prod(X-\alpha_i)$ в $\Om[X]$ однозначно). Значит,
$\sigma(E)=F[\sigma\alpha_1,\dots,\sigma\alpha_n]\subset E$, и $\beta'=\sigma(\beta)\in E$. Итак, все корни $g$ в $\Om$ лежат в $E$, т.\,е.\
$g$ раскладывается в $E[X]$ на линейные множители.
\end{proof}

\begin{keepfig}
\begin{center}
\begin{tikzcd}[column sep=large,row sep=large]
E=F[\beta][\alpha_1,\dots,\alpha_n]\arrow[r,cnew,dashed,"\sigma"] & \Om\supset E\\
F[\beta]\arrow[u,hook]\arrow[ur,cold,"\psi\colon\,\beta\,\mapsto\,\beta'"'] &
\end{tikzcd}
\end{center}
\figcap{К доказательству \rf{3.10a}: продолжение $\sigma$ переводит корни $\alpha_i$ многочлена $f$ в корни $f$, т.\,е.\ снова в
$\alpha_j$; поэтому $\sigma(E)\subset E$, и любой корень $\beta'$ многочлена $\mu_\beta$ лежит в $E$.}
\end{keepfig}

\begin{example}[3.10b: поле, порождённое одним корнем]
Пусть $E=F[\alpha]$, и минимальный многочлен $\mu_\alpha$ раскладывается в $E[X]$: $\mu_\alpha=\prod_{i=1}^n(X-\alpha_i)$
(корни перечислены с кратностями, $\alpha_1=\alpha$). По \rf{3.10a} $E/F$ нормально. Это можно увидеть и непосредственно,
получив попутно больше. Пусть $b\in E$; тогда $b=h(\alpha)$ для некоторого $h\in F[X]$. Положим
\[
r(X)=\prod_{i=1}^n\bigl(X-h(\alpha_i)\bigr).
\]
Коэффициенты $r$~--- симметрические многочлены от $\alpha_1,\dots,\alpha_n$ с коэффициентами из $F$, поэтому они выражаются
через коэффициенты $\mu_\alpha$ (основная теорема о симметрических многочленах, \rf{5.36}) и лежат в $F$. Так как
$r(b)=0$, минимальный многочлен $\mu_b$ делит $r$, и все корни $\mu_b$ находятся среди $h(\alpha_i)\in E$: $\mu_b$
раскладывается в $E[X]$. Более того:
\begin{enumerate}[label=\textup{(\roman*)}]
\item \emph{каждый корень $r$ является корнем $\mu_b$}: многочлен $\mu_b(h(X))\in F[X]$ обращается в нуль в $\alpha$, поэтому
делится на $\mu_\alpha$ и обращается в нуль во всех $\alpha_i$;
\item \emph{$r=\mu_b^k$, где $k=[F[\alpha]:F[b]]$.} Запишем $r=\mu_b^k\,d$ с $\mu_b\nmid d$. Если $d$ непостоянен, то его корень
является корнем $r$, а по (i)~--- корнем неприводимого $\mu_b$; тогда $\mu_b\mid d$~--- противоречие. Значит, $d=1$, и
$\deg\mu_\alpha=k\deg\mu_b$, т.\,е.\ $k=[F[\alpha]:F]/[F[b]:F]=[F[\alpha]:F[b]]$ (\rf{1.20}). В частности, $\deg\mu_b$ делит
$\deg\mu_\alpha$, и корень кратности $m$ многочлена $\mu_b$ встречается среди $h(\alpha_1),\dots,h(\alpha_n)$ ровно $km$ раз (если
$\mu_b$ сепарабелен~--- ровно $k$ раз).
\end{enumerate}
Нормальность $E$ нужна лишь для того, чтобы корни $\mu_b$ оказались в $E$. Если $\alpha$~--- любой алгебраический над $F$
элемент, $b=h(\alpha)$, а $\alpha_1,\dots,\alpha_n$~--- корни $\mu_\alpha$ (с кратностями) в каком-нибудь поле $\Om\supset F[\alpha]$,
где $\mu_\alpha$ раскладывается, то те же рассуждения (корни $d$ берутся в $\Om$) показывают, что
$r=\prod_i\bigl(X-h(\alpha_i)\bigr)\in F[X]$ и выполнены (i) и (ii): $r=\mu_b^{[F[\alpha]:F[b]]}$. Сепарабельность $\mu_\alpha$ при
этом нигде не используется. Например, для $\alpha=\sqrt[3]2$, $b=\alpha^2$ и $\Om=\CC$ имеем
$r=\prod_{j=0}^{2}\bigl(X-\omega^{2j}\sqrt[3]4\bigr)=X^3-4=\mu_b$ ($k=1$), хотя $\Q[\sqrt[3]2]/\Q$ не нормально.

Многочлен $r$~--- это характеристический многочлен умножения на $b$ в $F[\alpha]$ (\rf{5.44}; в сепарабельном случае ср.\ \rf{5.47b}). Пример: $\alpha=\sqrt2+\sqrt3$,
$\mu_\alpha=X^4-10X^2+1$, $b=\alpha^2=5+2\sqrt6$; тогда $r=(X^2-10X+1)^2$, $\mu_b=X^2-10X+1$, $k=2$.

Если $\mu_b$ сепарабелен и $k\cdot1_F\ne0$ (например, в характеристике $0$), то $\gcd(r,r')=\mu_b^{k-1}$ и
$\mu_b=r/\gcd(r,r')$. В характеристике $p$ эта формула может не работать: при $F=\F_2$, $E=\F_4=\F_2[\alpha]$,
$\alpha^2+\alpha+1=0$, $b=\alpha^3=1$ имеем $r=(X-1)^2=X^2+1$, $r'=0$, $\gcd(r,r')=r$, и $r/\gcd(r,r')=1\ne\mu_b=X-1$.
\end{example}

\begin{corollary}[3.11: теорема Артина]
Пусть $G$~--- конечная группа автоморфизмов поля $E$ и $F=E^G$. Тогда $E/F$~--- расширение Галуа с группой Галуа $G$, и
$[E:F]=(G:1)$.
\end{corollary}

\begin{proof}
То, что $E/F$~--- расширение Галуа, следует из \rf{3.10} ((c)$\Rightarrow$(d)); $\Gal(E/F)=\Aut(E/E^G)=G$ по \rf{3.5}; равенство
$[E:F]=|\Gal(E/F)|$ следует из \rf{3.2} (по \rf{3.10} поле $E$~--- поле разложения сепарабельного многочлена над $F$) или
прямо из доказательства \rf{3.5}.
\end{proof}

Итак, расширения Галуа~--- это в точности расширения вида $E/E^G$ с конечной группой $G\subset\Aut(E)$. Например, по
\rf{3.5c} расширения $k(X)/k(X+1/X)$ и $k(X)/k(j)$~--- расширения Галуа с группами порядков $2$ и $6$ (вторая изоморфна
$S_3$), а в характеристике $p$ расширение $k(X)/k(X^p-X)$~--- циклическое степени $p$. Так же по теореме Артина
$k(X_1,\dots,X_n)$~--- расширение Галуа поля симметрических рациональных функций с группой $S_n$ (\rf{5.38}--\rf{5.39}).

\begin{corollary}[3.12]
Каждое конечное сепарабельное расширение $E/F$ содержится в некотором расширении Галуа поля $F$.
\end{corollary}

\begin{proof}
Пусть $E=F[\alpha_1,\dots,\alpha_m]$, $f_i$~--- минимальный многочлен $\alpha_i$ над $F$ и $f$~--- произведение различных
многочленов среди $f_i$; он сепарабелен (пункт (c) после \rf{2.22}). Пусть $K$~--- поле разложения $f$ \emph{над $E$}
(\rf{2.7}). Оно порождено над $E$ корнями $f$, а $E$ порождено над $F$ корнями $f$ (элементами $\alpha_i$); значит, $K$
порождено корнями $f$ над $F$ и расщепляет $f$, т.\,е.\ $K$~--- поле разложения сепарабельного многочлена $f\in F[X]$ и по
\rf{3.10} является расширением Галуа поля $F$. По построению $K\supset E$.
\end{proof}

Поле разложения $f$ нужно строить над $E$, а не над $F$: абстрактно построенное поле разложения $f$ над $F$ содержит лишь
изоморфную копию $E$ (по \rf{2.12} в него есть $F$-гомоморфизм из $E$). Среди расширений Галуа, содержащих $E$, есть
наименьшее.

\begin{proposition}[3.12a: нормальное замыкание]
Пусть $E/F$~--- конечное расширение, $E=F[\alpha_1,\dots,\alpha_m]$, и $f$~--- произведение различных минимальных многочленов
$\alpha_i$ над $F$.
\begin{enumerate}[label=\textup{(\alph*)}]
\item Пусть $M\supset E$~--- нормальное расширение поля $F$. Подполе $L\subset M$, порождённое над $F$ корнями $f$ в $M$,
нормально над $F$, содержит $E$ и содержится в каждом подполе $M'\subset M$, содержащем $E$ и нормальном над $F$. Оно
порождено над $F$ всеми корнями в $M$ минимальных многочленов всех элементов $E$ и потому не зависит от выбора
образующих $\alpha_i$.
\item Поле разложения $L$ многочлена $f$ над $E$ нормально над $F$, и никакое собственное подполе $L$, содержащее $E$, не
нормально над $F$. Любые два таких расширения (нормальные над $F$, содержащие $E$ и минимальные с этим свойством)
$E$-изоморфны.
\item Если $E/F$ сепарабельно, то $L/F$~--- расширение Галуа.
\end{enumerate}
Поле $L$ из (b) называется \emph{нормальным замыканием} $E$ над $F$ (в сепарабельном случае~--- также \emph{замыканием
Галуа}); поле $L$ из (a)~--- нормальным замыканием $E$ в $M$.
\end{proposition}

\begin{proof}
(a) Каждый $f_i$ неприводим и имеет корень $\alpha_i\in M$, поэтому раскладывается в $M[X]$ ($M$ нормально); значит, $f$
расщепляется в $M$, и $L=F[\text{корни }f\text{ в }M]$~--- поле разложения $f$ над $F$, нормальное по \rf{3.10a}. Оно
содержит все $\alpha_i$, т.\,е.\ $E$. Если $M'\subset M$ содержит $E$ и нормально над $F$, то каждый $f_i$ раскладывается в
$M'[X]$; по однозначности разложения в $M[X]$ все корни $f_i$ в $M$ лежат в $M'$, и $L\subset M'$. Наконец, $L$ нормально и
содержит $E$, поэтому содержит все корни (в $M$) минимальных многочленов элементов $E$; обратно, корни $f$~--- среди них.

(b) Поле $L$ порождено над $E$ корнями $f$, а $E$ порождено над $F$ корнями $f$; значит, $L$~--- поле разложения $f$ и над
$F$, и по \rf{3.10a} нормально над $F$. Применяя (a) к $M=L$, видим, что каждое подполе $L$, содержащее $E$ и нормальное
над $F$, содержит $F[\text{корни }f]=L$. Если $L'$~--- другое нормальное над $F$ расширение $E$, минимальное с этим
свойством, то по (a) (для $M=L'$) $L'$ порождено над $F$, а значит, и над $E$ корнями $f$; т.\,е.\ $L'$~--- тоже поле
разложения $f$ над $E$, и по \rf{2.13} (над полем $E$) $L'\cong L$ над $E$.

(c) Если $E/F$ сепарабельно, то $f$ сепарабелен, и $L$~--- поле разложения сепарабельного многочлена (\rf{3.10}).
\end{proof}

Например, нормальное замыкание $\Q[\sqrt[3]2]$ над $\Q$~--- это $\Q[\sqrt[3]2,\omega]$ степени $6$, а нормальное замыкание
$\Q[\sqrt[4]2]$~--- поле $\Q[\sqrt[4]2,i]$ степени $8$ ($i\notin\Q[\sqrt[4]2]\subset\R$). Если $E/F$ уже нормально, то $L=E$.
На языке подгрупп нормальное замыкание описано в \rf{3.18}(b).

\begin{corollary}[3.13]
Пусть $F\subset M\subset E$. Если $E$~--- расширение Галуа поля $F$, то $E$~--- расширение Галуа поля $M$.
\end{corollary}

\begin{proof}
По \rf{3.10} $E$~--- поле разложения некоторого сепарабельного $f\in F[X]$. Оно же является полем разложения $f$,
рассматриваемого как многочлен из $M[X]$ (порождено над $M$ корнями $f$ и расщепляет $f$), а $f$ сепарабелен и над $M$.
Значит, $E/M$~--- расширение Галуа.
\end{proof}

\begin{example}[3.13a]
(1) Так же из \rf{3.10a}: если $E/F$ конечно и нормально, то нормально и $E/M$.

(2) Промежуточное поле $M$ не обязано быть нормальным над $F$: в башне $\Q\subset\Q[\sqrt[3]2]\subset\Q[\sqrt[3]2,\omega]$
верхнее поле~--- расширение Галуа поля $\Q$ (поле разложения $X^3-2$), а $\Q[\sqrt[3]2]/\Q$ не нормально (\rf{3.8}(a)).
Какие промежуточные поля нормальны над $F$, описывает основная теорема (\rf{3.17}(d)).

(3) \emph{Нормальность не транзитивна.} В башне $\Q\subset\Q[\sqrt2]\subset\Q[\sqrt[4]2]$ оба шага~--- расширения
Галуа степени $2$ ($\Q[\sqrt2]$~--- поле разложения $X^2-2$, $\Q[\sqrt[4]2]$~--- поле разложения $X^2-\sqrt2$ над
$\Q[\sqrt2]$, его корни $\pm\sqrt[4]2$), но $\Q[\sqrt[4]2]/\Q$ не нормально: $X^4-2$ имеет корень $\sqrt[4]2$, а корень
$i\sqrt[4]2$ не вещественен. Сепарабельность, напротив, транзитивна (\rf{3.15c}).
\end{example}

\begin{keepfig}
\begin{center}
\begin{tikzpicture}[every node/.style={font=\small}]
\node (Q) at (0,0) {$\Q$};
\node (Q2) at (0,1.4) {$\Q[\sqrt2]$};
\node (Q4) at (0,2.8) {$\Q[\sqrt[4]2]$};
\node[cnew] (Qi) at (0,4.2) {$\Q[\sqrt[4]2,\,i]$};
\draw (Q) -- node[right] {$2$, Галуа} (Q2);
\draw (Q2) -- node[right] {$2$, Галуа} (Q4);
\draw[cnew] (Q4) -- node[right] {$2$} (Qi);
\draw[decorate,decoration={brace,amplitude=5pt}] (-0.9,0.05) -- node[left=5pt,align=right] {$4$, не нормально} (-0.9,2.75);
\draw[cnew,decorate,decoration={brace,amplitude=5pt,mirror}] (2.3,0.05) -- node[right=5pt,align=left] {$8$, нормальное\\замыкание} (2.3,4.15);
\end{tikzpicture}
\end{center}
\figcap{Пример \rf{3.13a}(3): каждый шаг башни нормален, а вся башня~--- нет; наименьшее нормальное расширение $\Q$,
содержащее $\Q[\sqrt[4]2]$,~--- поле $\Q[\sqrt[4]2,i]$ степени $8$ (\rf{3.12a}).}
\end{keepfig}

\begin{remark}[3.14]
Пусть $E/F$~--- расширение Галуа с группой $G$, и $\alpha\in E$. Элементы $\alpha_1=\alpha,\alpha_2,\dots,\alpha_m$ орбиты $\alpha$ под
действием $G$ называются \emph{сопряжёнными} с $\alpha$. При доказательстве \rf{3.10} ((c)$\Rightarrow$(d), где $F=E^G$ по
\rf{3.10}(b)) мы показали, что минимальный многочлен $\alpha$ равен $\prod_{i=1}^m(X-\alpha_i)$. Иначе говоря,
\emph{сопряжённые с $\alpha$ элементы~--- это в точности корни его минимального многочлена} (все они лежат в $E$ и
попарно различны). В частности, число сопряжённых равно $[F[\alpha]:F]$.
\end{remark}

Если $E/F$ лишь нормально, это перестаёт быть верным: для $F=\F_p(T)$ и $E=F[\alpha]$, $\alpha^p=T$, группа $\Aut(E/F)$
тривиальна (\rf{3.3}(2)), орбита $\alpha$ состоит из одного элемента, а $\mu_\alpha=(X-\alpha)^p$ имеет степень $p$.

\begin{remark}[3.14a]
Пусть $E/F$~--- расширение Галуа с группой $G$, $\alpha\in E$, $\alpha_1,\dots,\alpha_n$~--- сопряжённые с $\alpha$ и $b=h(\alpha)$,
$h\in F[X]$. Формулу \rf{3.10b}(ii) теперь легко получить из теории групп. Отображение $\sigma\mapsto\sigma\alpha$ переводит $G$ на
орбиту $\alpha$, и каждый сопряжённый элемент имеет $|G_\alpha|$ прообразов, где $G_\alpha$~--- стабилизатор $\alpha$; так же для
$b$. Поскольку $\sigma b=h(\sigma\alpha)$,
\[
\Bigl(\prod_{i=1}^n\bigl(X-h(\alpha_i)\bigr)\Bigr)^{|G_\alpha|}=\prod_{\sigma\in G}(X-\sigma b)=\mu_b^{|G_b|},
\]
и по однозначности разложения $\prod_i(X-h(\alpha_i))=\mu_b^{k}$, где
$k=|G_b|/|G_\alpha|=(|G|/\deg\mu_b)/(|G|/\deg\mu_\alpha)=[F[\alpha]:F[b]]$ (\rf{3.14}). Итак, когда $\alpha$ пробегает свои
сопряжённые, $h(\alpha)$ пробегает сопряжённые с $b$, и каждый~--- ровно $[F[\alpha]:F[b]]$ раз.
\end{remark}

\begin{remark}[3.15]
Элемент $\alpha$ алгебраического расширения $E$ поля $F$ называется \emph{сепарабельным над $F$}, если его минимальный
многочлен над $F$ сепарабелен (\rf{2.22d}).

Доказательство \rf{3.12} показывает, что \emph{каждое конечное расширение $E=F[\alpha_1,\dots,\alpha_m]$, порождённое
сепарабельными элементами, сепарабельно} (никаких условий независимости на образующие не нужно: они могут быть
сопряжены и даже совпадать). Действительно, для произведения $f$ различных минимальных многочленов
$\alpha_i$ (а оно сепарабельно) поле разложения $K$ многочлена $f$ над $E$~--- расширение Галуа поля $F$ (\rf{3.10}). По
определению \rf{3.9} расширение $K/F$ сепарабельно, т.\,е.\ сепарабелен каждый элемент $K$, в частности каждый элемент
$E\subset K$. Например, если $\alpha$ сепарабелен, то сепарабельно всё простое расширение $F[\alpha]$: оно лежит в поле
разложения многочлена $\mu_\alpha$ над $F[\alpha]$, а это расширение Галуа. (Два доказательства, не использующие теорию
Галуа, даны в \rf{2.22g}.)

Поэтому элементы алгебраического расширения $E/F$, сепарабельные над $F$, образуют подполе $E_{\mathrm{sep}}$ поля $E$,
сепарабельное над $F$: если $\alpha,\beta\in E$ сепарабельны, то $F[\alpha,\beta]$ сепарабельно над $F$, и
$\alpha\pm\beta$, $\alpha\beta$, $\alpha^{-1}$ ($\alpha\ne0$) сепарабельны. Оно называется \emph{сепарабельным замыканием $F$ в $E$}.
Когда $E$ конечно над $F$, число $[E:F]_{\mathrm{sep}}=[E_{\mathrm{sep}}:F]$ называется \emph{сепарабельной степенью} $E$
над $F$.

Алгебраическое расширение $E/F$ называется \emph{чисто несепарабельным}, если единственные элементы $E$, сепарабельные
над $F$,~--- элементы $F$, т.\,е.\ $E_{\mathrm{sep}}=F$. Если $E$ конечно над $F$, то $E$ чисто несепарабельно над
$E_{\mathrm{sep}}$: пусть $\alpha\in E$ сепарабелен над $E_{\mathrm{sep}}$; тогда $E_{\mathrm{sep}}[\alpha]/E_{\mathrm{sep}}$ и
$E_{\mathrm{sep}}/F$~--- конечные сепарабельные расширения, по транзитивности (\rf{2.22h}, \rf{3.15c}) расширение
$E_{\mathrm{sep}}[\alpha]/F$ сепарабельно, и $\alpha\in E_{\mathrm{sep}}$. (С транзитивностью для алгебраических расширений,
\rf{3.15c}, то же верно для любого алгебраического $E/F$.)
\end{remark}

\begin{proposition}[3.15a: чисто несепарабельные расширения]
Пусть $\operatorname{char}F=p\ne0$ и $E/F$ алгебраично.
\begin{enumerate}[label=\textup{(\alph*)}]
\item Для каждого $\alpha\in E$ найдётся $e\ge0$, при котором $\alpha^{p^e}$ сепарабелен над $F$.
\item $E/F$ чисто несепарабельно $\iff$ для каждого $\alpha\in E$ некоторая степень $\alpha^{p^e}$ лежит в $F$ $\iff$ минимальный
многочлен каждого $\alpha\in E$ имеет вид $X^{p^e}-a$, $a\in F$.
\item Если $E/F$ конечно и чисто несепарабельно, то $[E:F]$~--- степень $p$.
\item Если $E/F$ конечно, то $[E:F]=[E:F]_{\mathrm{sep}}\cdot[E:E_{\mathrm{sep}}]$, и $[E:E_{\mathrm{sep}}]$~--- степень $p$
(\emph{несепарабельная степень}).
\end{enumerate}
\end{proposition}

\begin{proof}
(a) По упражнению 2-6 $\mu_\alpha=g(X^{p^e})$, где $g$ неприводим и сепарабелен; $\alpha^{p^e}$~--- корень $g$, т.\,е.\ сепарабелен.

(b) Если $E/F$ чисто несепарабельно, то элемент $\alpha^{p^e}$ из (a) сепарабелен и потому лежит в $F$. Пусть $\alpha^{p^e}\in F$ и
$e$ минимально; положим $a=\alpha^{p^e}$. При $e=0$ имеем $\mu_\alpha=X-a$. При $e\ge1$ элемент $a$ не является $p$-й
степенью в $F$: из $a=c^p$ следовало бы $(\alpha^{p^{e-1}})^p=c^p$, т.\,е.\ $\alpha^{p^{e-1}}=c\in F$ (Фробениус инъективен),
вопреки минимальности. По \rf{2.3a}(2) многочлен $X^{p^e}-a$ неприводим, т.\,е.\ это $\mu_\alpha$. Наконец, если
$\mu_\alpha=X^{p^e}-a=(X-\alpha)^{p^e}$ для каждого $\alpha$, то сепарабелен лишь случай $e=0$, т.\,е.\ $\alpha\in F$.

(c) Пусть $E=F[\alpha_1,\dots,\alpha_m]$ и $E_i=F[\alpha_1,\dots,\alpha_i]$. Каждый элемент $E$ имеет $p^e$-ю степень в $F\subset E_{i-1}$,
поэтому по (b) расширение $E_i/E_{i-1}$ чисто несепарабельно, и минимальный многочлен $\alpha_i$ над $E_{i-1}$ имеет вид
$X^{p^{e_i}}-a_i$. Значит, $[E_i:E_{i-1}]=p^{e_i}$, и $[E:F]$~--- степень $p$ (\rf{1.20}).

(d) Первое равенство~--- \rf{1.20}; второе утверждение следует из (c), так как $E/E_{\mathrm{sep}}$ чисто несепарабельно (\rf{3.15}).
\end{proof}

Например, пусть $p$ нечётно, $F=\F_p(T)$, и $E=F[\alpha]$, $\alpha^{2p}=T$ (многочлен $X^{2p}-T$ неприводим по Эйзенштейну
для простого элемента $T$ кольца $\F_p[T]$, \rf{1.16c}(c)). Элемент $\alpha^p$~--- корень сепарабельного многочлена
$X^2-T$, поэтому $F[\alpha^p]\subset E_{\mathrm{sep}}$. С другой стороны, для $\beta=\sum c_i\alpha^i$ ($c_i\in F[\alpha^p]$) имеем
$\beta^p=\sum c_i^p\alpha^{ip}\in F[\alpha^p]$, так что по (b) $E/F[\alpha^p]$ чисто несепарабельно, и элемент $E$, сепарабельный над $F$
(а значит, и над $F[\alpha^p]$, \rf{3.6a}), лежит в $F[\alpha^p]$. Итак, $E_{\mathrm{sep}}=F[\alpha^p]=F[\sqrt T]$,
$[E:F]_{\mathrm{sep}}=2$, $[E:E_{\mathrm{sep}}]=p$.

\subsection*{Критерии сепарабельности и транзитивность}

\begin{proposition}[3.15b]
Для конечного расширения $E/F$ следующие условия равносильны:
\begin{enumerate}[label=\textup{(\alph*)}]
\item $E/F$ сепарабельно;
\item $E$ порождено над $F$ конечным числом сепарабельных элементов;
\item $E$ содержится в некотором расширении Галуа поля $F$;
\item существует поле $\Om\supset F$, для которого число $F$-гомоморфизмов $E\to\Om$ равно $[E:F]$;
\item \textup{(}при $\operatorname{char}F=p\ne0$\textup{)} $E=F[E^p]$.
\end{enumerate}
\end{proposition}

\begin{proof}
(a)$\Rightarrow$(b): возьмём базис. (b)$\Rightarrow$(c)~--- это доказательство \rf{3.12}. (c)$\Rightarrow$(a): расширение
Галуа сепарабельно по определению, а вместе с ним и каждый его элемент.

(a)$\Rightarrow$(d). Пусть $E=F[\alpha_1,\dots,\alpha_m]$, $f$~--- произведение различных минимальных многочленов $\alpha_i$ (он
сепарабелен) и $\Om$~--- поле разложения $f$ над $E$. По \rf{2.12} $F$-гомоморфизмов $E\to\Om$ ровно $[E:F]$.

(d)$\Rightarrow$(a). Пусть $\beta\in E$, $g=\mu_\beta$, и $s$~--- число различных корней $g$ в $\Om$; $s\le\deg g$. Каждый
$F$-гомоморфизм $E\to\Om$ ограничивается до $F$-гомоморфизма $\psi\colon F[\beta]\to\Om$; таких $\psi$ ровно $s$ (\rf{2.1}(b)), и
каждый продолжается на $E$ не более чем $[E:F[\beta]]$ способами (\rf{2.14} для расширения $E/F[\beta]$, если считать $\Om$
расширением $F[\beta]$ посредством $\psi$). Поэтому
\[
[E:F]\le s\,[E:F[\beta]]\le\deg g\cdot[E:F[\beta]]=[E:F],
\]
откуда $s=\deg g$: у $g$ есть $\deg g$ различных корней, и он сепарабелен.

(a)$\Leftrightarrow$(e)~--- это \rf{2.22f}.
\end{proof}

Условие (d) удобно тем, что проверяется подсчётом. Оно даёт короткое доказательство транзитивности сепарабельности и
заодно отвечает на естественный вопрос: как выразить минимальный многочлен $\alpha$ над $F$ через его минимальный
многочлен над промежуточным полем.

\begin{proposition}[3.15c: транзитивность сепарабельности]
Пусть $F\subset E\subset K$~--- алгебраические расширения, $E/F$ конечно и сепарабельно, а $\alpha\in K$ сепарабелен над $E$.
Пусть $\Om\supset K$~--- поле, в котором раскладываются $\mu^F_\alpha$ и минимальные многочлены над $F$ образующих $E$
(например, поле разложения их произведения над $K$). Тогда
\begin{enumerate}[label=\textup{(\alph*)}]
\item $\alpha$ сепарабелен над $F$;
\item $\displaystyle\Phi(X):=\prod_{\sigma}(\sigma\mu^E_\alpha)(X)=\bigl(\mu^F_\alpha(X)\bigr)^{[E[\alpha]:F[\alpha]]}$, где $\sigma$ пробегает
$[E:F]$ различных $F$-гомоморфизмов $E\to\Om$, а $\sigma\mu^E_\alpha$~--- многочлен, полученный применением $\sigma$ к
коэффициентам $\mu^E_\alpha$.
\end{enumerate}
Следовательно, \emph{если $E/F$ и $K/E$~--- сепарабельные алгебраические расширения, то $K/F$ сепарабельно.}
\end{proposition}

\begin{proof}
Положим $L=E[\alpha]$, $d=[L:E]=\deg\mu^E_\alpha$. По \rf{2.12} (как в \rf{3.15b}, (a)$\Rightarrow$(d)) $F$-гомоморфизмов $E\to\Om$ ровно
$[E:F]$. Пусть $\sigma$~--- один из них. Так как $\mu^E_\alpha$ делит $\mu^F_\alpha$ в $E[X]$, то $\sigma\mu^E_\alpha$ делит
$\sigma\mu^F_\alpha=\mu^F_\alpha$ и потому раскладывается в $\Om$; он сепарабелен, как и $\mu^E_\alpha$ (изоморфизм
$E\to\sigma E$ сохраняет условие $\gcd(g,g')=1$). По \rf{2.4}(b) $\sigma$ продолжается до $L$ ровно $d$ способами $\tau$, и
значения $\tau(\alpha)$~--- это $d$ корней $\sigma\mu^E_\alpha$. Поэтому
\[
|\Hom_F(L,\Om)|=[E:F]\cdot d=[L:F],\qquad \Phi=\prod_\sigma\prod_{\tau|_E=\sigma}(X-\tau\alpha)=\prod_{\tau\in\Hom_F(L,\Om)}(X-\tau\alpha).
\]
Сгруппируем теперь $\tau$ по ограничениям $\psi=\tau|_{F[\alpha]}$. Различных $\psi$ столько же, сколько различных корней у
$\mu^F_\alpha$ в $\Om$ (\rf{2.1}(b)), т.\,е.\ не больше $\deg\mu^F_\alpha=[F[\alpha]:F]$, а каждый $\psi$ имеет не более
$[L:F[\alpha]]$ продолжений на $L$ (\rf{2.14}). Так как $[L:F]=[L:F[\alpha]]\,[F[\alpha]:F]$, в обеих оценках достигается
равенство: у $\mu^F_\alpha$ ровно $\deg\mu^F_\alpha$ различных корней (это (a)), и каждый $\psi$ имеет ровно $[L:F[\alpha]]$
продолжений. Значит, каждый корень $\beta=\psi(\alpha)$ многочлена $\mu^F_\alpha$ встречается среди $\tau(\alpha)$ ровно
$[L:F[\alpha]]$ раз, что и даёт (b).

Последнее утверждение. Пусть $\alpha\in K$ и $\mu^E_\alpha=X^n+b_{n-1}X^{n-1}+\dots+b_0$. Поле $E_0=F[b_0,\dots,b_{n-1}]\subset E$
конечно и порождено сепарабельными над $F$ элементами, поэтому сепарабельно над $F$ (\rf{3.15}). Многочлен
$\mu^E_\alpha\in E_0[X]$ неприводим над $E$, а значит, и над $E_0$; т.\,е.\ $\mu^{E_0}_\alpha=\mu^E_\alpha$, и $\alpha$ сепарабелен над
$E_0$. По (a) для $F\subset E_0\subset K$ элемент $\alpha$ сепарабелен над $F$.
\end{proof}

\begin{keepfig}
\begin{center}
\begin{tikzpicture}[every node/.style={font=\small},>=stealth]
\node (F) at (0,0) {$F$};
\node (Fa) at (-2.2,1.5) {$F[\alpha]$};
\node (E) at (2.2,1.5) {$E$};
\node (L) at (0,3) {$L=E[\alpha]$};
\node (Om) at (6.2,3) {$\Om$};
\draw[->] (F) -- (Fa);
\draw[->] (F) -- (E);
\draw[->] (Fa) -- (L);
\draw[->] (E) -- (L);
\draw[->,cnew,dashed] (L) -- node[below]{$\tau$: $[L:F]$ штук} (Om);
\draw[->,cold] (E) -- node[below right=-2pt]{$\sigma$: $[E:F]$ штук} (Om);
\draw[->,cnew] (Fa.north) .. controls (-2.2,4.7) and (4.2,4.7) .. node[above]{$\psi=\tau|_{F[\alpha]}$: не больше $\deg\mu^F_\alpha$ штук} (Om.north west);
\end{tikzpicture}
\end{center}
\figcap{Доказательство \rf{3.15c}: гомоморфизмы $L\to\Om$ считаются двумя способами~--- через $E$ (их $[E:F]\cdot d=[L:F]$) и через
$F[\alpha]$ (не больше $\deg\mu^F_\alpha\cdot[L:F[\alpha]]$). Равенство заставляет $\mu^F_\alpha$ иметь $\deg\mu^F_\alpha$ различных корней.}
\end{keepfig}

\begin{remark}[3.15d]
(1) Многочлен $\Phi$ из \rf{3.15c}~--- <<$\mu^E_\alpha$ с коэффициентами, заменёнными на сопряжённые>>~--- лежит в $F[X]$ и имеет
корень $\alpha$, а каждый его множитель $\sigma\mu^E_\alpha$ сепарабелен. Однако сам $\Phi$ сепарабельным быть не обязан, и
сепарабельность $\mu^F_\alpha$ не выводится из одного лишь того, что $\mu^F_\alpha\mid\Phi$; нужен подсчёт. Пример: $F=\Q$,
$E=\Q[\sqrt[3]2]$, $\alpha=\omega\sqrt[3]2$. Тогда $\mu^E_\alpha=X^2+\sqrt[3]2\,X+\sqrt[3]4$ (\rf{3.3}(1)), три гомоморфизма $E\to\CC$
переводят $\sqrt[3]2$ в $\theta\in\{\sqrt[3]2,\omega\sqrt[3]2,\omega^2\sqrt[3]2\}$, и
\[
\Phi=\prod_\theta\bigl(X^2+\theta X+\theta^2\bigr)=(X^3-2)^2,\qquad [E[\alpha]:\Q[\alpha]]=[\Q[\sqrt[3]2,\omega]:\Q[\omega\sqrt[3]2]]=2:
\]
каждый корень $X^3-2$ является корнем ровно двух множителей. Вообще, так как $\Phi=(\mu^F_\alpha)^k$, $k=[E[\alpha]:F[\alpha]]$, а
каждый множитель $\sigma\mu^E_\alpha$ имеет только простые корни, каждый сопряжённый с $\alpha$ элемент является корнем ровно $k$
множителей $\sigma\mu^E_\alpha$; $\Phi=\mu^F_\alpha$ тогда и только тогда, когда $E\subset F[\alpha]$, т.\,е.\ $k=1$.

(2) Транзитивность для конечных расширений другим способом (через критерий \rf{3.15b}(e)) доказана в \rf{2.22h}, а для
произвольных алгебраических расширений~--- ещё раз в \rf{6.8b}. Нормальность, в отличие от сепарабельности, не
транзитивна (\rf{3.13a}(3)).

(3) Из \rf{3.15c} следует, что для $F\subset E\subset K$ расширение $K/F$ сепарабельно тогда и только тогда, когда сепарабельны
$K/E$ и $E/F$ (необходимость~--- \rf{3.6a}).

(4) \emph{Коэффициенты нужно сопрягать согласованно.} Многочлен $\mu^F_\alpha$ не определяется многочленом $\mu^E_\alpha$ и
минимальными многочленами $\mu^F_{b_i}$ его коэффициентов $b_i$ по отдельности: в $\Phi$ все коэффициенты переводятся
\emph{одним и тем же} гомоморфизмом $\sigma$. Пример: $F=\Q$, $E=\Q[\sqrt2]$, $\alpha$ и $\alpha'$~--- корни многочленов
$X^2+\sqrt2X+\sqrt2$ и $X^2+\sqrt2X-\sqrt2$. У всех нестарших коэффициентов обоих многочленов минимальный многочлен над $\Q$
один и тот же, $X^2-2$, но по \rf{3.15c}(b)
\[
\Phi_\alpha=(X^2+\sqrt2X+\sqrt2)(X^2-\sqrt2X-\sqrt2)=X^4-2X^2-4X-2,\qquad \Phi_{\alpha'}=X^4-2X^2+4X-2 .
\]
Оба многочлена неприводимы (критерий Эйзенштейна при $p=2$), поэтому $\mu^\Q_\alpha=\Phi_\alpha\ne\Phi_{\alpha'}=\mu^\Q_{\alpha'}$.
Если же сопрягать коэффициенты независимо, получится произведение всех четырёх многочленов $X^2\pm\sqrt2X\pm\sqrt2$, т.\,е.\
$\Phi_\alpha\Phi_{\alpha'}$~--- одно и то же для $\alpha$ и $\alpha'$.

(5) \emph{Минимальный многочлен произведения.} Пусть $\alpha,\beta\ne0$~--- алгебраические над $F$ элементы некоторого
расширения, $n=\deg\mu_\alpha$, и $\Om\supset F[\alpha,\beta]$~--- поле, в котором раскладываются $\mu_\alpha$ и $\mu_\beta$. Многочлен
\[
f(X)=\prod_{\beta'}\beta'^{\,n}\mu_\alpha(X/\beta')=\prod_{\alpha',\beta'}(X-\alpha'\beta')
\]
($\alpha'$, $\beta'$ пробегают корни $\mu_\alpha$ и $\mu_\beta$ в $\Om$ с кратностями) лежит в $F[X]$: это
$\prod_{\beta'}g(X,\beta')$ для $g(X,Y)=Y^n\mu_\alpha(X/Y)\in F[X,Y]$, и его коэффициенты~--- симметрические многочлены от
корней $\beta'$ (\rf{5.36}). Так как $f(\alpha\beta)=0$, многочлен $\mu_{\alpha\beta}$ делит $f$, но $f$, вообще говоря, не
является степенью $\mu_{\alpha\beta}$~--- по той же причине, что в (4). Например, при $F=\Q$, $\alpha=1+\sqrt2$, $\beta=\sqrt[4]2$
\[
f=(X^4-8X^2-2)(X^4+8X^2-2),\qquad \mu_{\alpha\beta}=X^4-8X^2-2:
\]
любой гомоморфизм $\tau\colon\Q[\beta]\to\CC$ переводит $\alpha=1+\beta^2$ в $1+\tau(\beta)^2$, поэтому пары
$(1-\sqrt2,\pm\sqrt[4]2)$ и $(1+\sqrt2,\pm i\sqrt[4]2)$ ни из какого $\tau$ не получаются; они и дают лишний множитель
$X^4+8X^2-2$. Правильный многочлен даёт согласованное сопряжение: если $L=F[\alpha,\beta]$ сепарабельно над $F$, то
$F$-гомоморфизмов $L\to\Om$ ровно $[L:F]$ (\rf{2.12}), и вторая половина доказательства \rf{3.15c} (где использовалось только
это равенство), применённая к $L$ и элементу $\alpha\beta$ вместо $\alpha$, даёт
\[
\prod_{\tau\in\Hom_F(L,\Om)}\bigl(X-\tau(\alpha)\tau(\beta)\bigr)=\mu_{\alpha\beta}^{\,[L:F[\alpha\beta]]}.
\]
Так, при $\alpha=\sqrt2$, $\beta=\sqrt3$ обе конструкции дают $(X^2-6)^2=\mu_{\sqrt6}^2$, и $[\Q[\sqrt2,\sqrt3]:\Q[\sqrt6]]=2$.
То же верно для суммы $\alpha+\beta$ (с $\alpha'+\beta'$ вместо $\alpha'\beta'$ и $\prod_{\beta'}\mu_\alpha(X-\beta')$ вместо $f$).
\end{remark}

\begin{aside}[3.15e: сепарабельность и этальные алгебры]
Сепарабельность конечного расширения можно выразить через кольцо, которое получается при расширении скаляров. Пусть
$E=F[\alpha]\cong F[X]/(f)$, $f=\mu_\alpha$, и $\Om\supset F$~--- поле, в котором $f=\prod_j(X-\beta_j)^{m_j}$ с различными $\beta_j$.
Тогда $E\otimes_F\Om\cong\Om[X]/(f)$, и по китайской теореме об остатках
\[
E\otimes_F\Om\ \cong\ \prod_j\Om[X]/\bigl((X-\beta_j)^{m_j}\bigr).
\]
Если $f$ сепарабелен, все $m_j=1$ и $E\otimes_F\Om\cong\Om\times\dots\times\Om$ ($[E:F]$ сомножителей); если нет, то в
сомножителе с $m_j>1$ класс $X-\beta_j$~--- ненулевой нильпотент. Например, для $F=\F_p(T)$, $E=F[T^{1/p}]$ и $\Om=E$ получаем
$E\otimes_FE\cong E[X]/((X-T^{1/p})^p)$. Конечная $F$-алгебра $A$, для которой $A\otimes_F\Om\cong\Om^{\dim_FA}$ при некотором
поле $\Om\supset F$, называется \emph{этальной}. Итак, простое расширение сепарабельно тогда и только тогда, когда оно является
этальной алгеброй; по теореме о примитивном элементе (\rf{5.1}) это верно для всех конечных сепарабельных расширений, а для
несепарабельного $E$ кольцо $E\otimes_F\Om$ при $\Om$, расщепляющем минимальные многочлены, содержит ненулевые
нильпотенты (уже подкольцо $F[\beta]\otimes_F\Om$ для несепарабельного $\beta\in E$). Этальные алгебры~--- произведения конечных
сепарабельных расширений~--- ведут себя лучше полей (тензорное произведение полей обычно не поле, а тензорное
произведение этальных алгебр этально), и на их языке теория Галуа формулируется для них целиком; в этом курсе мы этим
языком не пользуемся.
\end{aside}

\begin{definition}[3.16]
Расширение $E/F$ называется \emph{циклическим} (соответственно \emph{абелевым}, \emph{разрешимым} и т.\,д.), если оно
является расширением Галуа с циклической (соответственно абелевой, разрешимой и т.\,д.) группой Галуа. Говоря об
\emph{экспоненте} абелева расширения, имеют в виду экспоненту его группы Галуа.
\end{definition}

Например, квадратичное расширение $E=F[\sqrt a]$ в характеристике $\ne2$ ($a\notin F^2$) циклично степени $2$: это поле
разложения сепарабельного многочлена $X^2-a$, и $\Gal(E/F)=\{\id,\sqrt a\mapsto-\sqrt a\}$. Расширение $\F_{p^n}/\F_p$
циклично, его группа порождена Фробениусом (\rf{4.20}); круговые поля $\Q[\zeta_n]$ абелевы (\rf{5.8}--\rf{5.10}); поле
разложения $X^3-2$ над $\Q$ разрешимо, но не абелево (его группа $S_3$); расширение $k(X)/k(X^p-X)$ из \rf{3.5c}(2)~---
циклическое степени $p$.
